\section{Rotation pairs and reversal}
Suppose $G=\gen{a,b}$, with $|a|,|b|\geq3$,
$(ab)^2=1$ and $\gen a\cap\gen b=1$.
If an involutory automorphism $\delta$ inverts both generators,
form $R=G\rtimes\gen\delta$ and put
\[
r_0=a\delta,\qquad r_1=\delta,\qquad r_2=\delta b.
\]
These are three distinct nonidentity involutions, the endpoints
commute, and they generate $R$. The adjacent subgroups are
$\gen a\rtimes\gen\delta$ and $\gen b\rtimes\gen\delta$;
their intersection is exactly $\gen\delta$ by the disjoint
cyclic subgroups in $G$. The complete rank-three intersection test
of Topic F proves the string C-group property. Thus $R$ is the
full regular group with the asserted type and order $2|G|$.
This gives the complete group-theoretic regular-doubling interface;
the classical string C-group--polytope correspondence is retained.


\begin{lemma}[Cyclic quotient criterion]\label{UF.CyclicQuotient}
Let $\phi:G\to H$ be a homomorphism and let $a,b\in G$.
If $\gen{\phi(a)}\cap\gen{\phi(b)}=1$ and the restriction of
$\phi$ to either $\gen a$ or $\gen b$ is injective, then
$\gen a\cap\gen b=1$.
\end{lemma}
\begin{proof}
Any common element maps into the trivial intersection downstairs.
It is therefore the identity by injectivity on the specified cyclic
subgroup. This is the quotient criterion used in source Proposition 2.1;
no finite-order or surjectivity hypothesis is needed for this step.
\end{proof}

\begin{lemma}[Abelian kernel reversal]\label{lem:abelian}
Suppose $G=P\rtimes\gen t$, where $P$ is abelian and $t^2=1$.
If $a\in P$ and $b=a^{-1}t$, an involutory automorphism of $G$
simultaneously inverts $a,b$.
\end{lemma}
\begin{proof}
Inversion on $P$ is an automorphism and commutes with the action
of $t$, so extends to an automorphism $\iota$ fixing $t$.
Follow it by conjugation with $a$. The result sends $a$ to $a^{-1}$
and $t$ to $a^{-1}ta$, hence sends $b=a^{-1}t$ to
$a(a^{-1}ta)=ta=b^{-1}$. Its square fixes $a$ and every element
of $P$, and fixes $t$ because $\iota(a)=a^{-1}$; thus it is
involutory. No presentation-order argument is needed.
\end{proof}

